% !TeX program = xelatex
\documentclass[11pt,a4paper]{ctexart}

\usepackage[a4paper,margin=2.35cm]{geometry}
\usepackage{amsmath,amssymb,mathtools,bm}
\usepackage{xcolor}
\usepackage[most]{tcolorbox}
\usepackage{enumitem,microtype,booktabs,hyperref}

\hypersetup{colorlinks=true,linkcolor=blue!50!black,urlcolor=blue!50!black}
\setlength{\parindent}{2em}
\setlength{\parskip}{0.35em}
\allowdisplaybreaks[3]

\definecolor{softyellow}{RGB}{255,249,225}
\definecolor{softgreen}{RGB}{238,249,241}
\newcommand{\diag}{\operatorname{diag}}
\newcommand{\e}{\mathrm{e}}
\newcommand{\vone}{\bm{1}}
\newcommand{\trans}{\mathsf{T}}

\newtcolorbox{checknote}[1][]{
  enhanced,breakable,colback=softyellow,colframe=orange!65!black,
  boxrule=0.7pt,arc=1.5mm,left=2mm,right=2mm,top=1.2mm,bottom=1.2mm,
  title=核对说明,#1
}
\newtcolorbox{resultbox}[1][]{
  enhanced,breakable,colback=softgreen,colframe=green!45!black,
  boxrule=0.7pt,arc=1.5mm,left=2mm,right=2mm,top=1.2mm,bottom=1.2mm,#1
}

\title{优化方法解 Newton 方程\\[0.4em]
\large Sinkhorn 候选方向与单参数最小残差}
\author{Basiliscus}
\date{2026年9月24日}

\begin{document}
\maketitle
\vspace{-2.2em}

\section*{符号与记法}
\begingroup
\small
\renewcommand{\arraystretch}{1.18}
\noindent
\begin{tabular*}{\textwidth}{@{}l@{\extracolsep{\fill}}l@{}}
\toprule
\textbf{符号} & \textbf{含义} \\
\midrule
$P,\ K$ & 当前传输矩阵与 Gibbs 核矩阵，$P=\diag(u)K\diag(v)$ \\
$a,\ b$ & 总质量相同且分量严格为正的目标边缘 \\
$\lambda,\ \mu$ & 对偶势；$u=\e^{-\lambda/\varepsilon},\ v=\e^{-\mu/\varepsilon}>0$ \\
$r_a=P\vone-a,\ r_b=P^{\trans}\vone-b$ & 当前两组边缘残差 \\
$m=P\vone,\ n=P^{\trans}\vone$ & 当前两组边缘质量向量 \\
$p,\ q$ & $u,v$ 的\emph{相对负增量}，即 $\Delta u=-u\odot p,\ \Delta v=-v\odot q$ \\
$\theta,\ h$ & 候选方向的平衡系数与外层时间步长 \\
$\vone,\ \odot,\ \oslash$ & 全一向量、逐元素乘法、逐元素除法 \\
$\diag(x),\ (\cdot)^{\trans}$ & 对角矩阵与转置；$\diag$ 只表示矩阵运算 \\
\bottomrule
\end{tabular*}
\par
\endgroup

\section{从 Newton 方程起步}

沿用《Sinkhorn 基础理解》中的缩放变量与符号约定。对边缘方程
$F(u,v)=(r_a,r_b)^{\trans}=0$，Newton 的\emph{实际增量}
$(\Delta u,\Delta v)$ 满足 $J_F(u,v)(\Delta u,\Delta v)^{\trans}=-F(u,v)$，即
\begin{equation}\label{eq:newton-uv}
 -\begin{pmatrix}
   \diag(Kv)&\diag(u)K\\
   \diag(v)K^{\trans}&\diag(K^{\trans}u)
  \end{pmatrix}
  \binom{\Delta u}{\Delta v}
 =\binom{r_a}{r_b}.
\end{equation}
这里 $\diag(v)K^{\trans}$ 是第二行的交叉块；它对 $u$ 求导，不能写成
$\diag(u)K^{\trans}$。在同一 $(u,v)$ 处定义
\begin{equation}\label{eq:pq}
 p=-\Delta u\oslash u,\qquad q=-\Delta v\oslash v,
 \qquad m=P\vone,\qquad n=P^{\trans}\vone.
\end{equation}
代入式~\eqref{eq:newton-uv}，得到与手稿中两行一一对应的方程
\begin{resultbox}
\begin{equation}\label{eq:block}
 \boxed{\begin{aligned}
   \diag(m)p+Pq&=r_a,\\
   P^{\trans}p+\diag(n)q&=r_b.
 \end{aligned}}
\end{equation}
\end{resultbox}
由于 $P=\diag(u)K\diag(v)$，例如
$\diag(Kv)\diag(u)=\diag(m)$、
$\diag(u)K\diag(v)=P$；第二行同理。

\begin{checknote}
手稿中在绝对量上分别近似 $\diag(m)p\approx m$ 与 $Pq\approx-a$，
相加形式上给出 $r_a=m-a$，但接近解时两项都可能是 $O(1)$，
需要以相反符号精确抵消。为避免这种消去，下面直接在\emph{残差尺度}
上分配 $r_a$。这里讨论的是在一个受限方向族内近似求解 Newton 方程，
不是要求这两项分别等于 $m$、$-a$。
\end{checknote}

\section{列边缘已归一时的 $\theta$ 方向}

先考虑刚完成一轮列归一化的状态：
\begin{equation}\label{eq:column-state}
 r_b=0,\qquad n=b,\qquad r_a=m-a.
\end{equation}
以冻结 $v$ 时的行 Sinkhorn 修正作为候选方向，令
\begin{equation}\label{eq:p-theta}
 p_S=\diag(m)^{-1}r_a=r_a\oslash m,
 \qquad p(\theta)=\theta p_S.
\end{equation}
这把手稿中「第一行两项承担多少」化为一个标量选择。
因为 $m>0$，式~\eqref{eq:p-theta} 有定义。若取 $\theta=1$ 且暂时\emph{冻结 $v$}，则
\begin{equation}\label{eq:row-exact}
 u^+=u- u\odot p_S
 =u\odot\bigl(a\oslash m\bigr)
 =a\oslash(Kv),
\end{equation}
这\emph{严格}是行 Sinkhorn 更新，无需近解近似。

下一步要求式~\eqref{eq:block} 的\emph{第二行线性方程}严格成立。
因为 $n=b$，直接取
\begin{equation}\label{eq:q-theta}
 q(\theta)=-\diag(b)^{-1}P^{\trans}p(\theta)
 =-\theta\bigl(P^{\trans}p_S\bigr)\oslash b.
\end{equation}
于是 $P^{\trans}p(\theta)+\diag(b)q(\theta)=0$。
不需要为 $q$ 另外解一个线性系统。

将式~\eqref{eq:p-theta}--\eqref{eq:q-theta} 代回第一行，定义
\begin{equation}\label{eq:s}
 s:=\diag(m)p_S-P\bigl[(P^{\trans}p_S)\oslash b\bigr]
 =r_a-P\bigl[(P^{\trans}(r_a\oslash m))\oslash b\bigr].
\end{equation}
第一行的 Newton 方程缺陷恰好是
\begin{equation}\label{eq:defect}
 \diag(m)p(\theta)+Pq(\theta)-r_a
 =\theta s-r_a.
\end{equation}
这也是手稿末尾「使第一行尽量精确成立」的一种可计算实现。

\section{把解方程化为一个小优化问题}

第一版先取欧氏范数，解一维最小残差问题：
\begin{resultbox}
\begin{equation}\label{eq:ls}
 \boxed{\theta_*\in\underset{\theta\in[0,\theta_{\max}]}{\operatorname{arg\,min}}
  \;\|\theta s-r_a\|_2^2,\qquad \theta_{\max}>0.}
\end{equation}
当 $\langle s,s\rangle>0$ 时，未加约束的解与有界解分别为
\begin{equation}\label{eq:theta}
 \theta_{\rm free}=\frac{\langle s,r_a\rangle}{\langle s,s\rangle},
 \qquad
 \theta_*=\min\!\bigl\{\theta_{\max},
          \max\{0,\theta_{\rm free}\}\bigr\}.
\end{equation}
\end{resultbox}
参数 $\theta_{\max}$ 只是对极大系数的保护，先固定而不自适应调节。
若 $s=0$ 且 $r_a\ne0$，则该一维方向无法减小线性方程缺陷；
该步标记为退化，改用标准 Sinkhorn 子步。若 $r_a=0$，
两组边缘残差都为零，直接停止。

定义相对 Newton 方程缺陷
\begin{equation}\label{eq:eta}
 \eta(\theta_*)=
 \frac{\bigl\|\binom{\diag(m)p(\theta_*)+Pq(\theta_*)-r_a}
                        {P^{\trans}p(\theta_*)+\diag(b)q(\theta_*)}\bigr\|_2}
      {\|\binom{r_a}{0}\|_2}
 =\frac{\|\theta_*s-r_a\|_2}{\|r_a\|_2}.
\end{equation}
它测量的是当前候选方向对\emph{线性 Newton 方程}的拟合程度；
仅凭 $\eta$ 较小，尚不能保证有限步迭代收敛。

\section{更新、有限步误差及第一轮实验}

按式~\eqref{eq:pq} 的符号，以步长 $h\in(0,1]$ 更新
\begin{equation}\label{eq:update}
 u^+=u\odot(\vone-hp(\theta_*)),\qquad
 v^+=v\odot(\vone-hq(\theta_*)).
\end{equation}
应确保两个括号逐分量为正；若不满足，缩短 $h$。
这与在 $\lambda,\mu$ 上作对数势增量并不等价。

为了精确区分「第二行线性方程成立」与「更新后列边缘可行」，
将 $P^+=\diag(\vone-hp)P\diag(\vone-hq)$ 代入得到
\begin{align}\label{eq:finite}
 r_a^+&=r_a-h\bigl[\diag(m)p+Pq\bigr]
                  +h^2 p\odot(Pq),\\
 r_b^+&=r_b-h\bigl[P^{\trans}p+\diag(n)q\bigr]
                  +h^2q\odot(P^{\trans}p).
\end{align}
在式~\eqref{eq:column-state}--\eqref{eq:q-theta} 下，特别有
\begin{equation}\label{eq:column-drift}
 r_b^+=h^2q\odot(P^{\trans}p)
 =-h^2\bigl(P^{\trans}p\bigr)\odot
                 \bigl[(P^{\trans}p)\oslash b\bigr].
\end{equation}
故第二行\emph{线性化}残差为零，但一般 $r_b^+\ne0$，
且此处的每个分量非正。为使下一轮仍从 $r_b=0$ 出发，
可在接受式~\eqref{eq:update} 后再做一次精确的列归一化
$v^+\leftarrow b\oslash(K^{\trans}u^+)$；这一额外操作应单独计入成本。

\begin{checknote}
$\theta=1$ 只保证\emph{冻结 $v$、仅执行 $u$ 子步}时得到式~\eqref{eq:row-exact}。
一旦同时执行 $q(1)$，更新就不是标准 Sinkhorn 一整轮；
$\theta=0$ 且 $r_b=0$ 时该构造不动。$\theta$ 也不必落在 $[0,1]$，
不能把它直接解释为「Sinkhorn 与完整 Newton 的插值比例」。
此外 $(p,q)=(c\vone,-c\vone)$ 是式~\eqref{eq:block} 的规范零方向，
完整 Newton 速度并非唯一；本文只用方程缺陷刻画近似程度。
\end{checknote}

\noindent\textbf{最小实验。}在每次列归一化后的状态，按
\eqref{eq:p-theta}、\eqref{eq:q-theta}、\eqref{eq:s}、\eqref{eq:theta}
计算方向和 $\theta_*$，记录 $\eta$，再执行有正性保障的式~\eqref{eq:update}
及下一次列归一化。第一轮仅比较标准 Sinkhorn 和此 $\theta$ 版本；
同时画 $\|r_a\|_1+\|r_b\|_1$ 随 $K/K^{\trans}$ 矩阵向量乘次数、
实际时间变化的曲线，并记录 $\theta_*$、$h$、$\eta$、
触发退化回退和缩步的次数。成本按实际的 $K/K^{\trans}$ 调用计，
包括额外列归一化。这样能先判断：小优化问题带来的方向信息，
是否抵得上这些额外乘法和同步更新的二次误差。

\section*{目前可确认的结论}

手稿提出的「第一行平衡、第二行尽量精确」可以严格落实为
式~\eqref{eq:ls}：第二行在线性层面精确，第一行在一维候选方向上
最小残差。它保留了 Sinkhorn 行子步的候选方向与 Newton 方程的块耦合，
但是否比相同计算预算下的 Sinkhorn 更有效，需要上述最小实验检验。

\end{document}
